\lesson{5}{Sep 27 2021 Mon (07:10:04)}{Operations on Complex Numbers}{Unit 1}

\subsubsection{Multiplying Complex Numbers}
\label{sub_sub_sec:multiplying_complex_numbers}

Multiplying complex numbers is a little bit different than adding and
subtracting complex numbers.

\begin{example}
  \label{exm:multiplying_complex_numbers}
  Let's simplify the following expression: $\sqrt{-12} \times \sqrt{-5}$

  \begin{equation}
    \label{eqt:multiplying_complex_numbers_1}
    \begin{split}
      \sqrt{-12} &\times \sqrt{-5} \\
      (\sqrt{-1} \times \sqrt{12}) &\times (\sqrt{-1} \times \sqrt{5}) \\
      i\sqrt{12} &\times i\sqrt{5}
    \end{split}
  .\end{equation}

  Once the imaginary numbers have been factored, multiply the imaginary numbers
  together and multiply the radical factors together.

  \begin{equation}
    \label{eqt:multiplying_complex_numbers_2}
    \begin{split}
      i\sqrt{12} &\times i\sqrt{5} \\
      i^{2}\sqrt{60} \\
      -2\sqrt{15}
    \end{split}
  .\end{equation}

  There is another method for multiplying complex numbers. You first multiply
  the square roots instead of factoring. Here's how that would look like:

  \begin{equation}
    \label{eqt:multiplying_complex_numbers_3}
    \begin{split}
      \sqrt{-12} &\times \sqrt{-5} \\
      \sqrt{60} \\
      \sqrt{4} &\times \sqrt{15} \\
      2\sqrt{15}
    \end{split}
  .\end{equation}

  You can see here (\ref{eqt:multiplying_complex_numbers_3}) that the negative
  was lost when the radicands were multiplied before factoring the imaginary
  number. This shows why the imaginary number must be factored first when
  multiplying radicals.
\end{example}

% subsubsection multiplying_complex_numbers (end)

\newpage
